% Supplemental material for the English main paper (egoanchor_en_ready_v4.tex).
% Compiles standalone: latexmk -g -xelatex -outdir=pdf egoanchor_en_supplementary.tex
% Covers the five items the main text promises:
%   S1 runtime parameters      <- main text Sec. 4 "Anchoring runtime"
%   S2 object-anchoring items  <- main text Sec. 5.3.3 "Object-anchoring items"
%   S3 standardized scales     <- main text Sec. 5.3.3 "User trust (TiA and S-TIAS)"
%   S4 complete statistics     <- main text Sec. 6.3
%   S5 scale reliability       <- main text Sec. 6.3.2
% Values are transcribed from: supplementary_runtime_parameters.tex (S1),
% material/EgoAnchor_Experiment3_Complete_Questionnaire_v5_3_Bilingual.md (S2, S3),
% tables/en/exp3_subjective.tex (S4), and
% data/experiments/experiment_3/analysis/results/experiment3_analysis.xlsx (S5).

\documentclass[review]{vgtc}

\usepackage{times}
\renewcommand*\ttdefault{txtt}
\usepackage{mathptmx}
\usepackage{makecell}

\usepackage{amsmath}
\usepackage{amssymb}
\usepackage{enumitem}
\hyphenation{Static-Lock Ego-Anchor}
\usepackage{array}
\usepackage{booktabs}
\usepackage{tabularx}
\usepackage{xcolor}

\onlineid{1118}
\vgtccategory{System}
\vgtcinsertpkg

% Number sections, tables and equations with an S prefix so that references
% cannot be confused with the main paper.
\renewcommand{\thesection}{S\arabic{section}}
\renewcommand{\thetable}{S\arabic{table}}
\renewcommand{\theequation}{S\arabic{equation}}

\title{EgoAnchor: Zero-Shot Dynamic Anchoring of Everyday Objects in Passthrough Mixed Reality --- Supplemental Material}

\author{}

\abstract{}

\begin{document}

\firstsection{Runtime Parameters}
\maketitle
\label{sec:supp-params}

This document contains the five items referenced in the main paper: runtime
parameters, the seven object-anchoring items, standardized-scale details, the
complete statistics for twelve subjective outcomes, and scale reliability.

Table~\ref{tab:supp-params} lists the complete parameter values of the anchoring
runtime. All parameters were fixed in advance and were not adjusted per scenario,
per object, or per condition across Experiments 1--3.

The locked-state gain is
$g_j = (1 - 2^{-\Delta t_j / h_{\mathrm{creep}}}) R_j (1 - \rho_j^{\mathrm{head}})$,
where $\rho_j^{\mathrm{head}}$ is the larger normalized head linear or angular speed,
clipped to $[0,1]$ using 0.3~m/s and 60~$^\circ$/s as full-tolerance scales. It
reduces the gain, while the motion criteria are relaxed by $1+3\rho_j^{\mathrm{head}}$
up to $4\times$; reliability criteria are unchanged. A 600~ms settling window
freezes motion-based release criteria after head motion stops.

The leaky cumulative sum used for the persistent-residual criterion is updated as
\begin{equation}
E_{p,j} = 2^{-\Delta t_j / h_E} E_{p,j-1} + \frac{\Delta t_j}{t_{\mathrm{ref}}} R_j \max\!\left(0,\, d_{p,j} - d_{\mathrm{db}}\right),
\label{eq:supp-cusum}
\end{equation}
and analogously for rotation. Release occurs when either channel exceeds its limit;
the observation consensus uses the same evidence half-life.

The comparison parameters were also fixed in advance: One-Euro uses position
cutoff/speed-coefficient/derivative-cutoff values of 0.8~Hz/6/2~Hz and rotation
values of 1~Hz/1/2~Hz; Smoothed KF Extrapolation uses a 0.18~s prediction horizon
and a 0.06~s correction half-life.

\paragraph{Rotation Kalman filter.}
Translation uses independent position--velocity constant-velocity filters on the
three Cartesian axes. Rotation is represented by three independent constant-velocity
filters on the local $\mathrm{SO}(3)$ tangent state $(\phi_i,\omega_i)$, where
$\phi_i$ is a rotation-vector component and $\omega_i$ is body-frame angular
velocity. The continuous process model uses white acceleration noise, with the
translational and rotational values in Table~\ref{tab:supp-params}; observations
are expressed as position and rotation-vector residuals. The first observation
initializes a normalized reference quaternion, zero rotation vector, and zero angular
velocity with an initial angular-velocity variance of $1.0~(\mathrm{rad}/\mathrm{s})^2$
per axis. Each subsequent quaternion is normalized, hemisphere-aligned with the
current estimate, and mapped to a shortest-arc relative logarithm. After correction,
the tangent origin is recentered at the updated orientation and the body-frame angular
velocity is transported by the $\mathrm{SO}(3)$ right Jacobian; the exponential map
recovers the output quaternion. These operations maintain a consistent quaternion
sign and tangent-space branch near $\pi$.

\begin{table*}[t]
\centering
\caption{Anchoring runtime parameters. The upper block covers observation admission
and history query, the middle block covers static locking (StaticLock), and the
lower block covers the validity rules and VCD scoring.}
\label{tab:supp-params}
\setlength{\tabcolsep}{6pt}
\begin{tabular}{@{}lll@{}}
\toprule
Parameter & Symbol & Value \\
\midrule
\multicolumn{3}{@{}l}{\emph{Observation admission and history query}} \\
Admission reliability threshold & $R_{\min}$ & 0.2 \\
Kalman white-acceleration noise (trans./rot.) & --- & 0.002~m$^2$/s$^3$ / 0.2~rad$^2$/s$^3$ \\
Kalman measurement noise (trans./rot.) & --- & $4\times10^{-6}$~m$^2$ / $4\times10^{-4}$~rad$^2$ \\
Query lag scale factor & $s$ & 1.15 \\
Minimum query lag & $\Delta_{\min}$ & 250~ms \\
State-age estimator coefficients (up/down) & --- & 0.5 / 0.05 \\
\midrule
\multicolumn{3}{@{}l}{\emph{Static locking (StaticLock)}} \\
Lock-entry linear speed threshold & $v_{\mathrm{th}}$ & 50~mm/s \\
Lock-entry angular speed threshold & $\omega_{\mathrm{th}}$ & 22~$^\circ$/s \\
Lock-entry reliability threshold & $R_{\mathrm{lock}}$ & 0.25 \\
Lock-entry dwell time & $\tau_{\mathrm{dwell}}$ & 350~ms \\
Lock-entry motion statistic smoothing & --- & 0.5 (per observation) \\
Deadband (trans./rot.) & $d_{\mathrm{db}}$ / $\theta_{\mathrm{db}}$ & 8~mm / 3~$^\circ$ \\
Locked-update gain half-life & $h_{\mathrm{creep}}$ & 2.7~s \\
CUSUM limit (trans./rot.) & $E_p^{\max}$ / $E_\theta^{\max}$ & 80~mm / 20~$^\circ$ \\
CUSUM evidence half-life & $h_E$ & 0.27~s \\
CUSUM reference interval & $t_{\mathrm{ref}}$ & 0.2~s \\
Accumulated drift limit (trans./rot.) & --- & 15~mm / 12~$^\circ$ \\
Speed-escape ratio and duration & --- & 2.5$\times$, 400~ms \\
Low-reliability release (threshold/duration) & --- & 0.3, 600~ms \\
Relock suppression duration & --- & 1.0~s \\
Seam transition per-frame decay & --- & 0.85 \\
Head settling window & --- & 600~ms \\
Head-motion full-tolerance scale (lin./ang.) & --- & 0.3~m/s / 60~$^\circ$/s \\
Head-motion tolerance amplification limit & --- & 4$\times$ \\
\midrule
\multicolumn{3}{@{}l}{\emph{Validity and VCD scoring}} \\
Tracking reliability floor & --- & 0.5 \\
Tracking-hold duration & --- & 450~ms without a reliable observation \\
Lost determination & --- & 1.0~s without a reliable observation \\
Reacquisition trigger (threshold/duration) & --- & 0.45, 600~ms \\
Reacquisition cooldown & --- & 3.0~s \\
Color/depth weights & $w_C$ / $w_D$ & 0.2 / 0.8 \\
Logarithmic floor & $\epsilon$ & 0.05 \\
Depth tolerance (floor/distance ratio) & --- & 5~mm / 0.02 \\
Inlier-rate weight and residual factor & --- & 0.6 / 2.5 \\
Valid depth coverage threshold & --- & 0.10 \\
Structural weight limit and IQR threshold & --- & 0.35 / 20~mm \\
LAB lightness channel weight & --- & 0.3 \\
\bottomrule
\end{tabular}
\end{table*}

\section{Experiment 1 Sampling and Confidence Intervals}\label{sec:supp-exp1-ci}

Table~\ref{tab:supp-exp1-sampling} gives the evaluation-segment counts, full formal-task
durations, and recorded reference-motion ranges for every controlled sequence family.
Each row is one continuous formal recording; $n$ is its number of marker-delimited
evaluation segments. Ranges include rest intervals and all recorded samples. A single
operator recorded all sequences; the intervals below therefore describe within-operator
repeatability.

\begin{table}[t]
\centering
\caption{Experiment~1 sampling. Duration is the full formal-task recording time;
reference-motion ranges are the full recorded ranges of instantaneous linear/angular
speeds across each task.}\label{tab:supp-exp1-sampling}
\begin{tabular}{@{}lcccc@{}}
\toprule
Sequence family & $n$ & \makecell{Duration\\(s)} & \makecell{Linear\\speed\\(m/s)} & \makecell{Angular\\speed\\($^\circ$/s)} \\
\midrule
Static head motion & 12 & 79.157 & 0--0 & 0--0 \\
Start--stop 6-DoF & 12 & 95.293 & 0--0.923 & 0--374.3 \\
Continuous translation & 16 & 77.257 & 0--3.453 & 0--385.9 \\
Continuous rotation & 12 & 65.062 & 0--0.722 & 0--903.7 \\
Occlusion recovery & 12 & 106.511 & 0--0 & 0--0 \\
\bottomrule
\end{tabular}
\end{table}

Table~\ref{tab:supp-exp1-ci} reports a bootstrap 95\% confidence interval for every
median in Tables~1--3 of the main paper. Each metric is aggregated within a
marker-delimited evaluation segment before the median across segments is taken, and
intervals come from 20{,}000 resamples of those per-segment values. Segment counts
are 12 for static registration, 16 (translation) and 12 (rotation) for dynamic
following, and 12 for occlusion recovery and motion response.

\begin{table*}[t]
\centering
\caption{Experiment~1 medians with bootstrap 95\% confidence intervals (20{,}000 resamples of
marker-delimited evaluation segments). Each metric is aggregated within a segment before the median across
segments is taken; values reproduce Tables~1--3 of the main paper. A single operator recorded all tasks.}
\label{tab:supp-exp1-ci}
\setlength{\tabcolsep}{4pt}
\begin{tabular}{@{}lcccc@{}}
\toprule
Metric & Arrival & Capture & One-Euro & EgoAnchor \\
\midrule
\multicolumn{5}{@{}l}{\emph{Static registration ($n=12$)}} \\
Head-motion coupling (mm) & 43.67 [35.87, 56.42] & 13.03 [11.40, 16.39] & 10.65 [7.99, 14.12] & 0.82 [0.58, 0.92] \\
Head-motion coupling ($^\circ$) & 6.49 [4.50, 8.48] & 5.40 [4.38, 8.19] & 3.29 [3.01, 4.40] & 0.33 [0.22, 0.37] \\
Registration error (mm) & 44.31 [37.10, 56.53] & 17.54 [13.89, 19.88] & 14.00 [11.95, 16.90] & 6.60 [5.18, 8.01] \\
Registration error ($^\circ$) & 9.65 [7.42, 11.35] & 8.26 [7.44, 10.50] & 6.96 [6.09, 8.17] & 4.39 [4.05, 5.02] \\
Static jitter (mm) & 11.70 [10.70, 16.88] & 4.30 [3.57, 7.01] & 0.85 [0.67, 1.57] & 0.06 [0.06, 0.08] \\
Static jitter ($^\circ$) & 2.98 [2.26, 3.77] & 2.60 [2.26, 3.11] & 0.40 [0.34, 0.50] & 0.03 [0.03, 0.03] \\
\midrule
\multicolumn{5}{@{}l}{\emph{Dynamic following, translation ($n=16$)}} \\
Effective latency (ms) & 202.5 [185.0, 217.5] & 255.0 [252.5, 262.5] & 380.0 [370.0, 385.0] & 360.0 [355.0, 365.0] \\
LA-RMSE (mm) & 18.57 [15.08, 20.81] & 17.14 [13.42, 18.29] & 15.69 [13.27, 20.35] & 9.12 [7.46, 11.27] \\
CT-RMSE (mm) & 84.74 [71.89, 92.70] & 103.63 [87.98, 110.61] & 117.61 [102.34, 130.03] & 125.83 [110.47, 137.82] \\
Residual jitter (mm) & 41.08 [34.13, 45.05] & 39.69 [33.87, 45.51] & 5.25 [4.27, 7.33] & 4.56 [4.10, 5.98] \\
\midrule
\multicolumn{5}{@{}l}{\emph{Dynamic following, rotation ($n=12$)}} \\
Effective latency (ms) & 255.0 [250.0, 262.5] & 257.5 [252.5, 265.0] & 385.0 [380.0, 390.0] & 345.0 [342.5, 350.0] \\
LA-RMSE ($^\circ$) & 6.09 [5.72, 6.53] & 5.84 [5.43, 6.04] & 5.55 [5.08, 6.15] & 4.81 [4.67, 5.29] \\
CT-RMSE ($^\circ$) & 25.44 [24.17, 30.74] & 25.35 [24.09, 31.16] & 34.10 [31.71, 41.76] & 31.88 [30.24, 39.20] \\
Residual jitter ($^\circ$) & 10.40 [9.64, 12.69] & 10.26 [9.66, 12.36] & 2.70 [2.56, 2.84] & 2.64 [2.39, 3.04] \\
\midrule
\multicolumn{5}{@{}l}{\emph{State-transition response ($n=12$)}} \\
Occlusion recovery (mm) & 18.23 [14.04, 36.83] & 16.93 [14.10, 36.95] & 10.41 [7.30, 14.59] & 4.85 [4.51, 13.93] \\
Occlusion recovery ($^\circ$) & 18.39 [11.77, 27.96] & 18.38 [11.90, 27.98] & 8.66 [5.95, 14.10] & 5.52 [4.88, 12.86] \\
Motion-response latency (ms) & 157.0 [134.9, 207.3] & 263.8 [229.8, 281.5] & 334.0 [326.7, 348.6] & 510.4 [446.7, 582.4] \\
\bottomrule
\end{tabular}
\end{table*}

\section{Object-Anchoring Items}\label{sec:supp-items}

The seven object-anchoring items were administered after each of the six
method--object blocks on a seven-point agreement scale (1 = strongly disagree,
7 = strongly agree). Q1--Q4 address phase-specific behavior; Q5--Q7 address
overall judgments.

\begin{enumerate}[label=Q\arabic*.]
  \item (Stationary stability) During static observation, the virtual content
        stayed in the same position without moving.
  \item (Motion attachment) While I picked up, moved, and rotated the object, the
        virtual content remained attached to the same location on the physical
        object.
  \item (Orientation consistency) While I rotated the object, the orientation of
        the virtual content stayed consistent with the orientation of the physical
        object.
  \item (Recovery consistency) After the occlusion was removed, the virtual content
        recovered to a position consistent with the physical object.
  \item (Positional correctness) The virtual content appeared at a position that
        matched where the physical object actually was.
  \item (Willingness to rely) I would be willing to rely on this anchoring method
        in an MR application that requires virtual content to be accurately
        attached to a physical object.
  \item (Stability--responsiveness balance) Overall, this method achieved an
        appropriate balance between anchor stability and responsiveness.
\end{enumerate}

\section{Standardized Scales}\label{sec:supp-scales}

AQ contains three-item Embedding Quality and Interaction Quality
subscales, administered after each block on seven-point scales. TiA contains the
six-item Reliability/Competence and four-item Understanding/Predictability
subscales, administered once per method on five-point scales; RC3, RC5, UP2, and
UP4 are reverse-scored as $6-\text{raw}$. S-TIAS contains three seven-point
intensity items and was administered once per method. The final questionnaire
recorded method preference, trust choice, preference strength, discrimination
confidence, and post-study discomfort as descriptive outcomes.

\section{Complete Subjective Statistics}\label{sec:supp-stats}

Table~\ref{tab:supp-stats} reports all twelve subjective outcomes. The seven custom
items and five standardized-scale outcomes form separate Holm--Bonferroni families.

% Supplement variant of tables/en/exp3_subjective.tex.
% Identical body; the caption resolves the two main-paper cross-references
% (\ref{sec:exp3-design}, \ref{fig:exp3-subjective}) to plain text so that the
% supplement compiles standalone. Numbers must stay in sync with the main table.
\begin{table*}[t]
\centering
\caption{Twelve within-participant subjective outcomes in Experiment 3, grouped according to the measurement structure described in Sec.~5.3 of the main paper. Object-anchoring items and AQ subscales are first averaged within condition across the three objects; TiA and S-TIAS are administered once per method. Median values are reported as One-Euro / EgoAnchor. $W$ and $p_{\mathrm{adj}}$ are the Wilcoxon signed-rank statistic and the within-family Holm-adjusted $p$ value, respectively. $r_{\mathrm{rb}}$ is the matched-pairs rank-biserial correlation, where positive values favor EgoAnchor. TiA uses five-point scales and all other outcomes use seven-point scales; arrows indicate that higher scores are more favorable. All outcomes contain 24 complete pairs, and zero differences are excluded from the Wilcoxon rank statistic. For outcomes marked with $\dagger$, all nonzero paired differences have the same sign and $r_{\mathrm{rb}}$ reaches its boundary value, so no bootstrap confidence interval is reported. Distributions are shown in Fig.~4 of the main paper.}
\label{tab:supp-stats}
\setlength{\tabcolsep}{2pt}
\begin{tabular}{@{}lcccc@{}}
\toprule
\makecell{Outcome} & \makecell{One-Euro /\\EgoAnchor\,$\uparrow$} & $W$ & $p_{\mathrm{adj}}$ & \makecell{$r_{\mathrm{rb}}$\\{[}95\% CI{]}} \\
\midrule
\multicolumn{5}{@{}l}{\textit{Phase-specific anchoring behavior}} \\
Stationary stability & 4.33 / 5.00 & 13.5 & $<.001$ & .90 [.69, 1.00] \\
Motion attachment & 5.00 / 5.17 & 67.5 & .267 & .29 [$-$.22, .77] \\
Orientation consistency & 4.67 / 5.00 & 81 & .162 & .41 [$-$.05, .78] \\
Recovery consistency & 4.17 / 5.00 & 6 & $<.001$ & .95 [.80, 1.00] \\
\addlinespace
\multicolumn{5}{@{}l}{\textit{Overall anchoring judgments}} \\
Positional correctness & 4.67 / 5.00 & 16 & .001 & .85 [.57, 1.00] \\
Willingness to rely & 4.00 / 5.00 & 16.5 & $<.001$ & .87 [.62, 1.00] \\
Stability-responsiveness balance & 4.50 / 4.83 & 6 & .003 & .90 [.65, 1.00] \\
\midrule
\multicolumn{5}{@{}l}{\textit{Augmentation quality}} \\
AQ Embedding Quality & 4.78 / 5.17 & 32 & .004 & .75 [.42, .96] \\
AQ Interaction Quality & 4.89 / 5.06 & 102.5 & .446 & .19 [$-$.29, .64] \\
\addlinespace
\multicolumn{5}{@{}l}{\textit{Trust}} \\
TiA Reliability/Competence & 4.17 / 4.58 & 0 & $<.001$ & 1.00$^{\dagger}$ \\
TiA Understanding/Predictability & 4.00 / 4.38 & 32 & .004 & .75 [.42, .96] \\
S-TIAS & 5.33 / 5.67 & 22.5 & .004 & .79 [.47, .97] \\
\bottomrule
\end{tabular}
\end{table*}

\section{Scale Reliability}\label{sec:supp-reliability}

\textbf{Reliability coefficients.} Internal consistency by scale and method
($N=24$). Cronbach's $\alpha$ and McDonald's $\omega$ describe the present sample.
AQ coefficients use scores averaged across the three objects; TiA and S-TIAS were
administered once per method.
\begin{tabular}{@{}llccc@{}}
\toprule
Scale & Method & Items & $\alpha$ & $\omega$ \\
\midrule
\multicolumn{5}{@{}l}{\emph{Averaged across three objects}} \\
AQ Embedding Quality & One-Euro & 3 & .768 & .773 \\
AQ Embedding Quality & EgoAnchor & 3 & .769 & .770 \\
AQ Interaction Quality & One-Euro & 3 & .504 & .525 \\
AQ Interaction Quality & EgoAnchor & 3 & .892 & .899 \\
\midrule
\multicolumn{5}{@{}l}{\emph{Single method-level administration}} \\
TiA Reliability/Competence & One-Euro & 6 & .792 & .796 \\
TiA Reliability/Competence & EgoAnchor & 6 & .751 & .777 \\
TiA Understanding/Predictability & One-Euro & 4 & .565 & .580 \\
TiA Understanding/Predictability & EgoAnchor & 4 & .769 & .756 \\
S-TIAS & One-Euro & 3 & .672 & .698 \\
S-TIAS & EgoAnchor & 3 & .697 & .740 \\
\bottomrule
\end{tabular}

\end{document}
